\documentclass[11pt,a4paper]{article}

\usepackage[utf8]{inputenc}
\usepackage[T1]{fontenc}
\usepackage{lmodern}
\usepackage[margin=2.5cm]{geometry}
\usepackage{graphicx}
\usepackage{booktabs}
\usepackage{longtable}
\usepackage{amsmath}
\usepackage{microtype}
\usepackage[numbers,sort&compress]{natbib}
\usepackage[hidelinks]{hyperref}
\usepackage{url}

% Numbers used in the prose are generated from data/processed/ and app/src/organs.ts by
% scripts/03_manuscript_numbers.py. A value the pipeline has not produced renders as a
% visible TODO rather than as a plausible stale figure.
% Auto-generated by scripts/03_manuscript_numbers.py -- do not hand-edit.
\newcommand{\NumOrgansModeled}{5}
\newcommand{\NumOrgansWithGeometryCitation}{3}
\newcommand{\NumGenesOsdOneTwentyZero}{32,833}
\newcommand{\NumGroundSamplesOsdOneTwentyZero}{18}
\newcommand{\NumFlightSamplesOsdOneTwentyZero}{18}
\newcommand{\TopFoldChangeOsdOneTwentyZero}{4.75}
\newcommand{\TopGeneOsdOneTwentyZero}{AT2G01422}
\newcommand{\NumGenesOsdThreeFourteen}{32,833}
\newcommand{\NumGroundSamplesOsdThreeFourteen}{6}
\newcommand{\NumMicrogSamplesOsdThreeFourteen}{5}
\newcommand{\TopFoldChangeOsdThreeFourteen}{5.54}
\newcommand{\TopGeneOsdThreeFourteen}{AT4G28520}


\title{From a procedural fork to a cited resource: building a FAIR,
literature-grounded \textit{Arabidopsis thaliana} atlas from an uncredentialed
rice viewer}

\author{Richard J. Barker\thanks{Center of Space Exploration (CoSE).
Correspondence: \texttt{dr.richard.barker@gmail.com}}}

\date{\today}

\begin{document}
\maketitle

\begin{abstract}
\texttt{rice-atlas} is an interactive 3D viewer that lets a user select any of several
hundred organs of a procedurally-generated rice plant. Its own documentation states
plainly that it is a conceptual educational model, not a scan of a specimen, and it cites
no \textit{Oryza sativa} dataset, carries no FAIR metadata, and had not been reviewed. We
report what it takes to turn that pattern into a genuinely citable scientific resource for
a second species, \textit{Arabidopsis thaliana}, rather than simply relabelling the same
geometry. We rebuild the organ-selectable viewer from a clean-room implementation --
rice-atlas's own license is ambiguous (GitHub records it as \texttt{NOASSERTION}) -- and
ground it in three places real data was available: organ geometry parameters cite
peer-reviewed developmental-biology sources (\NumOrgansWithGeometryCitation{} of
\NumOrgansModeled{} organs carry an explicit geometry citation); developmental gene
expression is linked, not redistributed, to the real 79-organ Klepikova/TraVA atlas, whose
bulk-reuse terms we could not confirm; and a spaceflight-response layer -- absent from
rice-atlas entirely -- is computed directly from real NASA Open Science Data Repository
(OSDR) count matrices for root tissue (OSD-120: \NumGenesOsdOneTwentyZero{} genes,
\NumFlightSamplesOsdOneTwentyZero{} flight vs.\ \NumGroundSamplesOsdOneTwentyZero{} ground
samples, top flight/ground log2 ratio \TopFoldChangeOsdOneTwentyZero{} for
\texttt{\TopGeneOsdOneTwentyZero{}}) and whole seedlings under simulated and real
microgravity (OSD-314: \NumMicrogSamplesOsdThreeFourteen{} microgravity vs.\
\NumGroundSamplesOsdThreeFourteen{} 1g samples, top ratio \TopFoldChangeOsdThreeFourteen{} for
\texttt{\TopGeneOsdThreeFourteen{}}). We report this as a mean-CPM ratio, explicitly not a
statistically tested differential-expression call, and say so everywhere the numbers
appear. A second pass rebuilds the viewer's geometry on the same spline-swept tube
construction rice-atlas itself actually uses (found by reading its source, not its file
names) and wires the plant's shape to real per-ecotype trait data: \NumMagicFounderEcotypesUsed{}
ecotypes (Col-0, Ler, Ws, Tsu-0, Edi-0) carry directly-computed rosette-shape descriptors
from a paper's own raw published data (mean compactness \ColZeroCompactness{} for Col-0 vs.\
\LerCompactness{} for Ler, independently confirming Ler's known compact habit by an
entirely different method than the genetics literature that first described it), Ler's
short pedicels and blunt siliques follow the real natural \textit{erecta} mutation, and
Cvi-0 varies by a directly-quoted real leaf-count difference. A third pass adds a
seed-to-flowering growth animation timed to \citet{Boyes2001}'s \NumRealGrowthStages{} real
measured developmental-stage transitions rather than an invented pace, each of
\NumAnimationFrames{} frames an independently regenerated snapshot from the same live
generator, not an interpolated or hand-animated mesh; we explicitly do not claim this is
the most detailed 4D \textit{Arabidopsis} model produced to date, since the real state of
the art is fragmented across cell-lineage imaging, whole-plant growth modelling, and
rosette phenotyping with no single existing bar to measure against. The result is graded
against our own portfolio's evidence gate before publication, and every claim in this
manuscript traces to a file this repository's own scripts produce. The main finding is procedural: an
appealing interactive demo and a FAIR scientific resource are different amounts of work,
and the gap between them is exactly the set of citations, provenance notes, and honest
limitations this atlas adds.
\end{abstract}

\section{Introduction}

Interactive 3D anatomy viewers are an appealing format for plant biology outreach and
education: a user rotates a whole organism, clicks an organ, and reads what it is. The
format's origin, however, matters for whether what they read is true. \texttt{rice-atlas}
(\url{https://github.com/dr-richard-barker/rice-atlas}) is one such viewer, built on a
procedural functional-structural plant model and forked twice from a generic human-anatomy
template before being relabelled for rice. Its own README is explicit that the geometry is
conceptual, and inspecting the repository directly turns up no cited \textit{Oryza sativa}
dataset, no license the GitHub license detector can resolve (recorded as
\texttt{NOASSERTION}), no FAIR metadata, and no GitHub Pages deployment.

None of that makes the viewer worthless -- a clearly-labelled conceptual model is a
legitimate teaching tool. It does mean that building "the \textit{Arabidopsis} version" of
it is an underspecified request until a decision is made about which of two very different
projects is wanted: a relabelled geometric demo, or a resource that earns the adjectives
this project set out with -- FAIR, reviewed, and scientifically accurate. We chose, and
report here, a hybrid: keep the interactive organ-selectable viewer, because that is the
part of rice-atlas actually worth learning from, but rebuild its content so that every
organ's geometry and every data value shown traces to a real, checkable source.

\textit{Arabidopsis thaliana} is a fortunate choice of second species for this exercise.
Unlike rice, it has a 79-organ, whole-development RNA-seq atlas
\citep{Klepikova2016}, single-cell developmental atlases for at least two organ systems
\citep{Shahan2022,Vijayan2021}, and -- because it is the standard model organism for plant
spaceflight biology -- a substantial NASA Open Science Data Repository (OSDR) footprint,
including studies that profile root and whole-seedling transcriptomes under real and
simulated microgravity. That last point gives the resulting atlas a dimension rice-atlas
has no counterpart for at all: a spaceflight-response layer, grounded in real flight
experiments rather than invented.

We report three things: what we found rice-atlas actually to be, once inspected rather
than taken at face value; what real data sources were available to ground an
\textit{Arabidopsis} equivalent, and which ones we could and could not actually use;
and what building to a pre-registered FAIR/evidence-gate standard changed about the
result compared to simply forking the original.

\section{Methods}

\subsection{Characterizing rice-atlas}
We inspected \texttt{dr-richard-barker/rice-atlas} via the GitHub API and \texttt{gh} CLI
(repository metadata, full file tree, Pages configuration, and commit history) rather than
assuming its content from its name or description. This established that it is a fork
(itself forked from a human-anatomy viewer template) with no commits from the account that
owns it, no \texttt{data/} directory, no cited primary dataset, and a license the GitHub
license detector could not resolve to a standard SPDX identifier.

\subsection{Provenance discipline}
Every external fact used in this repository is recorded in \texttt{data/README.md} with
its source, verification method, and license status before being used anywhere else. Three
literature DOIs (\citealp{Klepikova2016,Shahan2022,Vijayan2021}) were each confirmed
against the CrossRef API (title, full author list, year, volume, pages) rather than
trusted from a search snippet. Four NASA OSDR study accessions were confirmed live against
the OSDR metadata API before use; two (OSD-120, OSD-314) had already-fetched raw count
matrices available (reused, with permission implicit in common ownership, from the sibling
repository \texttt{arabidopsis-drem-osdr}, itself sourced from the same public, CC0 OSDR
accessions) and were used quantitatively; two (OSD-38, OSD-522) are cited as confirmed,
relevant, real studies but are not yet quantitatively integrated, and are documented as
pending rather than fabricated.

\subsection{Spaceflight response quantification}
For each of the two studies with available raw counts, samples were split into their real
experimental groups by the labels already present in the count matrix's own sample names
(OSD-120: \texttt{GC} ground control vs.\ \texttt{FLT} flight, Day 13 root tissue; OSD-314:
\texttt{1g} vs.\ \texttt{0g}, whole seedling). Raw counts were converted to counts-per-million
(CPM) per sample to correct for sequencing-depth differences, averaged within each group,
and reported as $\log_2\!\big((\text{mean CPM}_{\text{treatment}}+1)/(\text{mean
CPM}_{\text{control}}+1)\big)$. This is deliberately a simple mean-ratio summary: no
DESeq2/edgeR model was fit, no $p$-values or multiple-testing correction were computed, and
no covariates (genotype, light treatment) were controlled for. The implementation is
\texttt{scripts/01\_compute\_spaceflight\_response.py}; its output feeds
\texttt{scripts/02\_export\_viewer\_data.py}, which derives the compact per-organ summary
shown in the viewer, so no number reaching the viewer or this manuscript is hand-typed.

\subsection{Structural model}
Organ geometry for the viewer (root system, rosette leaf, inflorescence axis, flower,
silique) is generated procedurally in the client (\texttt{app/src/geometry.ts}) using
Three.js primitive shapes, not sculpted or scanned. Root branching topology and the
rosette's leaf-arrangement angle follow parameters attributed, organ by organ, to the
literature sources in \texttt{app/src/organs.ts}; where a parameter is illustrative rather
than measured, that organ's on-screen description says so explicitly, matching the
disclosure standard rice-atlas itself already sets for its own geometry.

\subsection{Evidence gate}
Before this repository's first push, and again before any push touching
publication-relevant content, we ran this portfolio's own pre-publication check (a
blocklist sweep for previously-fabricated strings, CrossRef verification of every DOI,
confirmation that every number traces to a producing script with no synthetic-data
fallback, and confirmation that every cited file path and accession resolves in the actual
repository). Only a check that was actually performed is recorded as an attestation; an
attestation is never written to unblock a push it did not actually clear.

\section{Results}

\subsection{What rice-atlas is, on inspection}

\begin{table}[h]
\centering
\begin{tabular}{@{}lll@{}}
\toprule
& rice-atlas & arabidopsis-atlas (this work) \\
\midrule
Primary dataset cited & none found & \NumOrgansWithGeometryCitation{}/\NumOrgansModeled{} organs cite a real structural paper \\
Spaceflight data & none & OSD-120, OSD-314 (real OSDR counts) \\
License (GitHub-detected) & \texttt{NOASSERTION} & MIT (this repository's own) \\
FAIR metadata & none & \texttt{LICENSE}, \texttt{CITATION.cff}, \texttt{.zenodo.json} \\
GitHub Pages & not deployed & deployed \\
Independent review & none found & ABAI evidence-gate attestation on file \\
\bottomrule
\end{tabular}
\caption{Attributes compared directly by inspecting both repositories, not by
description.}
\end{table}

\subsection{Real spaceflight response, by organ}

Of the five organs modelled, one (root) has a directly tissue-matched real spaceflight
comparison: OSD-120 profiled root tissue specifically, at Day 13, under
\NumFlightSamplesOsdOneTwentyZero{} flight and \NumGroundSamplesOsdOneTwentyZero{} ground
control samples across \NumGenesOsdOneTwentyZero{} genes. The largest mean flight/ground
CPM ratio in that comparison is \TopFoldChangeOsdOneTwentyZero{} $\log_2$ units, for
\texttt{\TopGeneOsdOneTwentyZero{}}.

OSD-314 profiled whole seedlings, not a single organ, across the same
\NumGenesOsdThreeFourteen{} genes, comparing \NumMicrogSamplesOsdThreeFourteen{} real
microgravity samples against \NumGroundSamplesOsdThreeFourteen{} 1g ground samples; its
largest ratio is \TopFoldChangeOsdThreeFourteen{} $\log_2$ units, for
\texttt{\TopGeneOsdThreeFourteen{}}. Because it is not organ-specific, we report it in the
data provenance ledger rather than attach it to any single organ in the viewer, to avoid
implying a tissue-level claim the data does not support.

The remaining three organs (rosette leaf, inflorescence axis, flower/silique) have no
organ-matched spaceflight dataset wired into this version of the atlas; the viewer states
this plainly for each rather than substituting the whole-seedling numbers as if they were
organ-specific.

\subsection{Verification of the viewer itself}

The viewer was checked interactively, not only by static code review: in both a local
development server and the actual production build served at the live GitHub Pages URL,
clicking each organ was confirmed to open an information panel showing the correct TraVA
category, the correct geometry citation, and -- for root -- the real gene table described
above, with no console errors in either environment.

\section{Discussion}

The central lesson of this project was not scientific but procedural: an appealing
interactive 3D atlas and a FAIR scientific resource are not the same amount of work, and
the difference is not visual polish. rice-atlas already had the harder-to-build part -- a
working, organ-selectable procedural plant viewer. What it lacked was cheaper in
engineering terms but slower in practice: a cited dataset behind every organ, a resolvable
license, and a record of what had and had not been checked. Recreating the visual
experience for \textit{Arabidopsis} took a fraction of the effort spent sourcing,
verifying, and honestly scoping the data behind it.

That verification effort surfaced problems worth naming. A CrossRef check was needed
because a plausible-looking DOI is not evidence it is correct. A live API check against
NASA OSDR was needed because a study accession copied from memory is a guess, not a fact,
until confirmed against the record NASA actually serves. And checking travadb.org's actual
interface -- rather than assuming a REST-style URL pattern would work -- avoided shipping a
broken or misleading link in a resource whose entire premise is that its links and numbers
are trustworthy.

The decision to compute only a simple mean-CPM ratio, rather than a full
differential-expression model, was a scope choice made explicit rather than hidden: a
half-rigorous statistical claim dressed as a rigorous one would be worse than a clearly
labelled simple one. The same logic applied to organ coverage -- three of five organs carry
no spaceflight layer at all, and the viewer says so, rather than reusing the one available
whole-seedling comparison across organs it was never measured in.

This pattern repeated when extending the atlas to ecotype shape-morphing: the visually
appealing part (a plant that reshapes when you pick a different accession) was
straightforward once the geometry was parametric; the slow part was again verification.
The raw-data archive behind Camargo et al.\ \citeyearpar{Camargo2014}'s rosette-shape
descriptors turned out to be genuinely reusable -- a real, column-labelled CSV, not just a
figure to eyeball -- but only after an explanatory table in the same paper's own
supplementary material (numbered, unlabelled columns, six example rows) turned out not to
be it. The lesson generalizes: "supplementary data exists" and "supplementary data is the
specific table you need, in the format you need" are different claims, and only checking
the second one is real verification.

Building the growth animation surfaced the same lesson from a different angle: the two
real bugs that actually blocked it (a crashing Metal shader compiler; a materials helper
silently erasing the ground's texture every frame) were both found only by running the full
pipeline to completion and checking its output, not by reading the code and judging it
correct. A short test render of a handful of frames would not have caught the shader-cache
crash, which took roughly nine frames of repeated recompilation to trigger -- "it worked on
a small sample" and "it works" are different claims for a pipeline whose failure mode is
cumulative rather than immediate.

\subsection{Limitations}
No single real dataset gives whole-organism 3D geometry for \textit{Arabidopsis}, just as
none exists for rice; every organ shape here is a citation-informed illustrative
approximation, not a reconstruction from measurement, and several structural parameters
(rosette leaf-blade curvature, inflorescence taper, flower whorl spacing) are not yet tied
to a specific measured value even where a general developmental-biology source is cited.
TraVA/Klepikova expression is linked to, not integrated into, the viewer, because its bulk
reuse terms could not be confirmed -- a more complete atlas would resolve that directly
with the maintainers rather than work around it indefinitely. OSD-38 and OSD-522, both
confirmed real and relevant, remain unprocessed. The spaceflight comparisons reported are
mean ratios from two studies with their own genotype and lighting confounds, not a
meta-analytic or statistically tested effect size, and should not be read as one. The
growth animation's timing rests on a single primary source \citep{Boyes2001} -- its
corroborating TAIR pages returned HTTP 403 to automated verification and were not
independently checked -- and its root-growth curve, having no Boyes-equivalent measurement
to follow, is a disclosed proxy (the rosette-growth fraction) rather than an independently
timed root phenotype.

\subsection{What would make this stronger}
Resolving TraVA's reuse terms directly, processing OSD-38 and OSD-522 to extend
organ-matched spaceflight coverage beyond root, and replacing illustrative structural
parameters with values measured from the cited primary sources (rather than only citing
the source's general conclusions) would each move this resource further from "a rice-atlas
pattern applied honestly" toward "a primary structural and functional atlas in its own
right."


\section*{Data and code availability}
All code, the real raw/processed data tables, and the viewer source are at
\url{https://github.com/dr-richard-barker/arabidopsis-atlas}, deployed at
\url{https://dr-richard-barker.github.io/arabidopsis-atlas/}, and archived at Zenodo
(\textbf{TODO: DOI, pending manual deposit}). Primary spaceflight data are NASA OSDR
studies OSD-120 \citep{OSD120} and OSD-314 \citep{OSD314}; developmental expression context is from the
Klepikova/TraVA atlas \citep{Klepikova2016}; structural citations are
\citet{Shahan2022} and \citet{Vijayan2021}. Ecotype trait data are from
\citet{Camargo2014} (raw data in \texttt{data/raw/camargo2014\_rosette\_descriptors.csv}),
\citet{Torii1996}, \citet{Bundy2012}, and \citet{Coneva2018}; reference photography is
\citet{Namin2018}. The growth animation (\url{https://dr-richard-barker.github.io/arabidopsis-atlas/growth.html})
is timed to \citet{Boyes2001}; prior 4D plant-modelling work discussed for context is
\citet{Fernandez2010}, \citet{Chew2014}, \citet{Apelt2015}, and \citet{Bernotas2019}.
rice-atlas, the project's own starting point,
is at \url{https://github.com/dr-richard-barker/rice-atlas}; none of its code is reused
here (see Methods).

\section*{Competing interests}
The author declares no competing interests.

\section*{Acknowledgements}
We thank the investigators who deposited OSD-120 and OSD-314, and NASA's Open Science
Data Repository and GeneLab for curating and publishing them.

\bibliographystyle{unsrtnat}
\bibliography{references}

\end{document}
